\section{Related Work}
\label{sec:related}

% Prose v1 (2026-06-10). Toutes les cles \cite ci-dessous existent et sont
% verifiees dans references.bib. OWNER (Genest et al.) cite en 3e personne.
% Cles encore manquantes (PDF non fourni), laissees en commentaire si besoin :
%   BERT (Devlin 2019), survey DL-NER (Li et al. 2022), InstructUIE, GPT-NER.

OPENER sits at the intersection of three lines of research: open-world and zero-shot NER, pretrained entity representations refined by contrastive learning, and energy-aware machine learning.
We review each in turn and state how OPENER differs.

\subsection{Closed-set and Span-based NER}
\label{sec:related:closed}
Traditional NER is cast as supervised sequence labeling over a fixed inventory of entity types, typically with a pretrained encoder followed by a token or span-level classifier.
Early neural taggers replaced hand-made features with recurrent encoders over a CRF layer: the BiLSTM-CRF \cite{lample2016ner} paired character and word-level representations, and the BiLSTM-CNN-CRF \cite{ma2016end} added a character-level CNN to obtain a fully end-to-end sequence labeler.
The pretrained encoder paradigm subsequently increased in-domain accuracy further, fine-tuning a bidirectional transformer such as BERT \cite{devlin2019bert} with a lightweight classification head.
Beyond flat tagging, span-based formulations enumerate candidate spans and classify them directly, which simplifies the handling of nested and discontinuous entities \cite{nguyen2023span}: biaffine NER \cite{yu2020biaffine} scores every span with a biaffine layer borrowed from dependency parsing, while an MRC framework \cite{li2020mrc} recasts typing as machine reading comprehension, querying the sentence with a natural-language description of each type, a use of label semantics our zero-shot head also exploits.
These models reach high accuracy on in-domain benchmarks such as CoNLL-2003 \cite{tjongkimsang2003conll}, as documented in a comprehensive survey of deep learning for NER \cite{li2022survey}, but their label set is frozen at training time.
Adding a new entity type requires re-annotating data and retraining the model, which is the central limitation that open-world NER seeks to remove.

\subsection{Open-world and Zero-shot NER}
\label{sec:related:openworld}
Open-world NER aims to recognise entity types that were not fixed at training time.
GLiNER \cite{zaratiana2024gliner} treats the candidate type names as queries to a bidirectional encoder, a DeBERTaV3 backbone \cite{he2023debertav3}, and scores every span against them in a single forward pass, \emph{via a span-type similarity matrix}, offering a fast alternative to autoregressive tagging.
GNER \cite{ding2024gner} instead casts NER as text generation with a sequence-to-sequence model \cite{raffel2020t5}: it regenerates the input sentence with inline type tags (e.g. ``\texttt{Marie Curie(person) discovered radium(chemical)}''), and shows that explicitly modeling the negative, non-entity tokens in this output sharpens type boundaries.
Close to our work, OWNER \cite{genest2025owner} is an unsupervised open-world pipeline that decouples mention detection from entity typing: it embeds detected mentions with an encoder refined by contrastive learning and then clusters them into dynamically named types.
Zero-shot typing has also been driven by the semantics of the type names themselves: one line of work leverages natural-language type \emph{descriptions} \cite{aly2021zeroshot} so that classes unseen at training time can still be recognised, a premise our label-name prototypes share.
Progress in the low-resource regime is in turn anchored by dedicated benchmarks, notably Few-NERD \cite{ding2021fewnerd}, a large human-annotated dataset organised as a two-level hierarchy of $8$ coarse and $66$ fine-grained types.
Other approaches require a general-purpose large language model to perform extraction directly, as in ChatIE \cite{wei2024chatie}, or focus on recovering entities that detectors tend to miss \cite{cai2025handling}.
OPENER adopts the decoupled detect-then-type view of OWNER, but replaces the trained encoder and the unsupervised clustering with a pretrained, Matryoshka-truncatable embedding space and a balanced linear classifier, so that no encoder is trained from scratch.

\subsection{LLM-based and Few-shot NER}
\label{sec:related:llm}
A complementary line distils or instruction-tunes large language models for open NER.
UniversalNER \cite{zhou2024universalner} distils a chat model into a smaller LLaMA-based student \cite{touvron2023llama} over tens of thousands of entity types, while GoLLIE \cite{sainz2024gollie} conditions a code language model on annotation guidelines to follow unseen schemas.
Few-shot methods reduce supervision further through prompt optimization \cite{tang2024fsponer}, knowledge-graph-guided contrastive learning \cite{zhang2024kcl}, or token-level contrastive objectives such as CONTaiNER \cite{das2022container}.
These methods are accurate but rest on models with billions of parameters, whose inference cost, latency, and energy use are substantial, and which may not even fit on commodity hardware.
At the small end of this family, however, a quantized instruction-tuned model does fit a commodity GPU: as a concrete, runnable reference point, we include Qwen2.5-1.5B-Instruct \cite{qwen2025} in 4-bit, prompted to return entities as JSON, and report where it lands relative to the dedicated open-world systems (\autoref{sec:exp:results}).
This cost is precisely what motivates the lightweight design and the energy-aware evaluation of OPENER.

\subsection{Entity Representations and Contrastive Learning}
\label{sec:related:repr}
OPENER relies on a pretrained sentence embedder rather than a task-specific encoder.
Nomic Embed \cite{nussbaum2024nomic} provides an open, long-context model trained with the Matryoshka objective \cite{kusupati2022matryoshka}, whose representations can be truncated to lower dimensions without retraining, trading quality for compute.
Such embedders are based on the Sentence-BERT line of work \cite{reimers2019sbert}, and we sharpen entity-type geometry with a triplet margin objective in the spirit of FaceNet \cite{schroff2015facenet}.
Refining a contrastive objective by mining and re-weighting hard negatives is a recurring recipe, as in GCSE \cite{lai2025gcse}, which down-weights false hard negatives with a Gaussian decay, close in spirit to our error-driven hard-negative mining.
The choice of base encoder is made empirically: we compare several strong open sentence encoders (E5 \cite{wang2022e5}, BGE \cite{xiao2024cpack}, MPNet \cite{song2020mpnet} and mxbai \cite{li2024aoe}) on entity typing under a common protocol (\autoref{tab:ablation}). Nomic v1.5 \cite{nussbaum2024nomic} edges them out while additionally offering native Matryoshka truncation, so we adopt it as the base and fine-tune it.
Embedding labels and entities in a shared space \cite{wang2018joint}, modeling entity types with Gaussian components \cite{zhang2025contrastivegaussian}, generating synthetic entity-bearing text for augmentation \cite{hu2023entity}, and using mixtures to augment training data \cite{abbahaddou2025gmmgnn} are recurring ideas in low-resource NER, all grounded in classical Gaussian mixture models fitted by expectation-maximization \cite{reynolds2009gmm, dempster1977em}.
Unlike methods that train a dedicated encoder, OPENER takes transfer learning in its purest form: it keeps a general-purpose pretrained space largely frozen and only applies a light contrastive fine-tuning, testing whether such a representation is already discriminative enough for entity typing.

\subsection{Efficient and Green AI}
\label{sec:related:green}
A growing body of work argues that efficiency should be reported alongside accuracy.
Prior work has quantified the financial and environmental cost of training NLP models \cite{strubell2019energy}, promoted efficiency as a first-class evaluation criterion under the \emph{Green AI} agenda \cite{schwartz2020greenai}, and analysed the carbon footprint of machine-learning training and how to reduce it \cite{patterson2022carbon}.
Complementary efforts argue for the systematic reporting of energy and carbon and provide a framework to track them per run \cite{henderson2020systematic}, and take a whole-lifecycle, system-level view of AI's environmental footprint across data, algorithms, and hardware \cite{wu2022sustainable}.
Most relevant to deployment, a study of \emph{inference}-time cost \cite{luccioni2024power} finds general-purpose generative models to be orders of magnitude more expensive per query than task-specific ones, precisely the gap a lightweight pipeline like ours seeks to close.
Tools such as CodeCarbon \cite{luccioni2019quantifying} make per-run energy and carbon estimates practical.
However, open-world NER systems are still compared almost exclusively on quality.
Instead, OPENER reports a three-axis comparison of quality, measured by Adjusted Mutual Information \cite{vinh2010ami}, inference speed, and energy, placing the frugality of each method on an equal footing with its accuracy.

In summary, existing open-world NER methods either train a dedicated encoder \cite{genest2025owner} or query costly large language models \cite{zhou2024universalner, sainz2024gollie}, and are compared on accuracy alone.
OPENER fills both gaps: it reuses pretrained components with only a light contrastive step and a balanced linear classifier, and it is evaluated jointly on quality, speed and energy.
